% 本文件由 tools/build_book.py 从 dist/reports/assertions.json 生成，请勿手改。
\begin{longtable}{@{}p{0.06\textwidth}p{0.32\textwidth}p{0.22\textwidth}p{0.16\textwidth}@{}}
\caption{章节 5.2 数值断言（生成时间 2026-10-19T17:44:22-0400）}\\
\toprule
类别 & 断言 & 理论值 & 实测值 \\
\midrule
\endfirsthead
\toprule
类别 & 断言 & 理论值 & 实测值 \\
\midrule
\endhead
\bottomrule
\endlastfoot

B & lcurve\_\allowbreak{}tradeoff/log10\_\allowbreak{}lam 全有限 & 0 个非有限值 & 0 \checkmark \\
B & lcurve\_\allowbreak{}tradeoff/log10\_\allowbreak{}x\_\allowbreak{}norm 全有限 & 0 个非有限值 & 0 \checkmark \\
B & lcurve\_\allowbreak{}tradeoff/log10\_\allowbreak{}residual 全有限 & 0 个非有限值 & 0 \checkmark \\
B & lcurve\_\allowbreak{}tradeoff/error\_\allowbreak{}curve 全有限 & 0 个非有限值 & 0 \checkmark \\
B & lcurve\_\allowbreak{}tradeoff/curvature 全有限 & 0 个非有限值 & 0 \checkmark \\
B & lcurve\_\allowbreak{}tradeoff/dof 全有限 & 0 个非有限值 & 0 \checkmark \\
B & lcurve\_\allowbreak{}tradeoff/gcv 全有限 & 0 个非有限值 & 0 \checkmark \\
B & lcurve\_\allowbreak{}tradeoff/criterion\_\allowbreak{}log10\_\allowbreak{}lam 全有限 & 0 个非有限值 & 0 \checkmark \\
B & lcurve\_\allowbreak{}tradeoff/criterion\_\allowbreak{}error 全有限 & 0 个非有限值 & 0 \checkmark \\
B & lcurve\_\allowbreak{}tradeoff/criterion\_\allowbreak{}valid 全有限 & 0 个非有限值 & 0 \checkmark \\
B & lcurve\_\allowbreak{}tradeoff/corner\_\allowbreak{}log10\_\allowbreak{}lam 全有限 & 0 个非有限值 & 0 \checkmark \\
B & lcurve\_\allowbreak{}tradeoff/picard\_\allowbreak{}sigma 全有限 & 0 个非有限值 & 0 \checkmark \\
B & lcurve\_\allowbreak{}tradeoff/picard\_\allowbreak{}coeff 全有限 & 0 个非有限值 & 0 \checkmark \\
B & lcurve\_\allowbreak{}tradeoff/picard\_\allowbreak{}ratio 全有限 & 0 个非有限值 & 0 \checkmark \\
B & lcurve\_\allowbreak{}tradeoff/picard\_\allowbreak{}below\_\allowbreak{}noise 全有限 & 0 个非有限值 & 0 \checkmark \\
B & lcurve\_\allowbreak{}tradeoff/solution\_\allowbreak{}true 全有限 & 0 个非有限值 & 0 \checkmark \\
B & lcurve\_\allowbreak{}tradeoff/solution\_\allowbreak{}lam 全有限 & 0 个非有限值 & 0 \checkmark \\
B & filter\_\allowbreak{}shapes/index 全有限 & 0 个非有限值 & 0 \checkmark \\
B & filter\_\allowbreak{}shapes/log10\_\allowbreak{}sigma 全有限 & 0 个非有限值 & 0 \checkmark \\
B & filter\_\allowbreak{}shapes/filter\_\allowbreak{}tsvd 全有限 & 0 个非有限值 & 0 \checkmark \\
B & filter\_\allowbreak{}shapes/filter\_\allowbreak{}tikhonov 全有限 & 0 个非有限值 & 0 \checkmark \\
B & filter\_\allowbreak{}shapes/filter\_\allowbreak{}landweber 全有限 & 0 个非有限值 & 0 \checkmark \\
B & filter\_\allowbreak{}shapes/filter\_\allowbreak{}landweber\_\allowbreak{}valid 全有限 & 0 个非有限值 & 0 \checkmark \\
B & filter\_\allowbreak{}shapes/filter\_\allowbreak{}tv\_\allowbreak{}probe 全有限 & 0 个非有限值 & 0 \checkmark \\
B & filter\_\allowbreak{}shapes/filter\_\allowbreak{}tv\_\allowbreak{}probe\_\allowbreak{}mixed 全有限 & 0 个非有限值 & 0 \checkmark \\
B & filter\_\allowbreak{}shapes/dof\_\allowbreak{}tsvd\_\allowbreak{}curve 全有限 & 0 个非有限值 & 0 \checkmark \\
B & filter\_\allowbreak{}shapes/error\_\allowbreak{}tsvd\_\allowbreak{}curve 全有限 & 0 个非有限值 & 0 \checkmark \\
B & filter\_\allowbreak{}shapes/dof\_\allowbreak{}tikhonov\_\allowbreak{}curve 全有限 & 0 个非有限值 & 0 \checkmark \\
B & filter\_\allowbreak{}shapes/error\_\allowbreak{}tikhonov\_\allowbreak{}curve 全有限 & 0 个非有限值 & 0 \checkmark \\
B & filter\_\allowbreak{}shapes/log10\_\allowbreak{}lam\_\allowbreak{}curve 全有限 & 0 个非有限值 & 0 \checkmark \\
B & filter\_\allowbreak{}shapes/dof\_\allowbreak{}landweber\_\allowbreak{}curve 全有限 & 0 个非有限值 & 0 \checkmark \\
B & filter\_\allowbreak{}shapes/error\_\allowbreak{}landweber\_\allowbreak{}curve 全有限 & 0 个非有限值 & 0 \checkmark \\
B & filter\_\allowbreak{}shapes/error\_\allowbreak{}matched 全有限 & 0 个非有限值 & 0 \checkmark \\
B & image\_\allowbreak{}deconvolution/image\_\allowbreak{}true 全有限 & 0 个非有限值 & 0 \checkmark \\
B & image\_\allowbreak{}deconvolution/image\_\allowbreak{}data 全有限 & 0 个非有限值 & 0 \checkmark \\
B & image\_\allowbreak{}deconvolution/image\_\allowbreak{}recon 全有限 & 0 个非有限值 & 0 \checkmark \\
B & image\_\allowbreak{}deconvolution/image\_\allowbreak{}abs\_\allowbreak{}error 全有限 & 0 个非有限值 & 0 \checkmark \\
B & image\_\allowbreak{}deconvolution/image\_\allowbreak{}residual 全有限 & 0 个非有限值 & 0 \checkmark \\
B & image\_\allowbreak{}deconvolution/history\_\allowbreak{}k 全有限 & 0 个非有限值 & 0 \checkmark \\
B & image\_\allowbreak{}deconvolution/error\_\allowbreak{}history 全有限 & 0 个非有限值 & 0 \checkmark \\
B & image\_\allowbreak{}deconvolution/residual\_\allowbreak{}history 全有限 & 0 个非有限值 & 0 \checkmark \\
B & image\_\allowbreak{}deconvolution/log10\_\allowbreak{}sigma 全有限 & 0 个非有限值 & 0 \checkmark \\
B & image\_\allowbreak{}deconvolution/log10\_\allowbreak{}lam\_\allowbreak{}sweep 全有限 & 0 个非有限值 & 0 \checkmark \\
B & image\_\allowbreak{}deconvolution/error\_\allowbreak{}lam\_\allowbreak{}sweep 全有限 & 0 个非有限值 & 0 \checkmark \\
B & bias\_\allowbreak{}variance/log10\_\allowbreak{}lam 全有限 & 0 个非有限值 & 0 \checkmark \\
B & bias\_\allowbreak{}variance/bias2\_\allowbreak{}analytic 全有限 & 0 个非有限值 & 0 \checkmark \\
B & bias\_\allowbreak{}variance/variance\_\allowbreak{}analytic 全有限 & 0 个非有限值 & 0 \checkmark \\
B & bias\_\allowbreak{}variance/mse\_\allowbreak{}analytic 全有限 & 0 个非有限值 & 0 \checkmark \\
B & bias\_\allowbreak{}variance/bias2\_\allowbreak{}mc 全有限 & 0 个非有限值 & 0 \checkmark \\
B & bias\_\allowbreak{}variance/variance\_\allowbreak{}mc 全有限 & 0 个非有限值 & 0 \checkmark \\
B & bias\_\allowbreak{}variance/mse\_\allowbreak{}mc 全有限 & 0 个非有限值 & 0 \checkmark \\
B & bias\_\allowbreak{}variance/mse\_\allowbreak{}mc\_\allowbreak{}std 全有限 & 0 个非有限值 & 0 \checkmark \\
B & bias\_\allowbreak{}variance/error\_\allowbreak{}rel\_\allowbreak{}analytic 全有限 & 0 个非有限值 & 0 \checkmark \\
B & bias\_\allowbreak{}variance/error\_\allowbreak{}rel\_\allowbreak{}mc 全有限 & 0 个非有限值 & 0 \checkmark \\
B & bias\_\allowbreak{}variance/dof 全有限 & 0 个非有限值 & 0 \checkmark \\
B & bias\_\allowbreak{}variance/local\_\allowbreak{}slope 全有限 & 0 个非有限值 & 0 \checkmark \\
B & bias\_\allowbreak{}variance/local\_\allowbreak{}slope\_\allowbreak{}variance 全有限 & 0 个非有限值 & 0 \checkmark \\
A & 四个可视化的数组数 / 构建耗时 / JSON 合法性 & pack\_frame 要求 meta 只有 number/\allowbreak{}bool/\allowbreak{}null、\allowbreak{}数组全有限、\allowbreak{}字符串走 strings & lcurve\_tradeoff: 17 数组 /\allowbreak{} 44 meta /\allowbreak{} 0.01s；filter\_shapes: 16 数组 /\allowbreak{} 30 meta /\allowbreak{} 0.04s；image\_deconvolution: 11 数组 /\allowbreak{} 27 meta /\allowbreak{} 0.23s；bias\_variance: 13 数组 /\allowbreak{} 27 meta /\allowbreak{} 0.04s(\allowbreak{}最慢 0.23s) \checkmark \\
A & 参数角点 + t 端点的帧数与最慢单帧耗时 / 非有限 meta 数 & 奇异状态必须变成可显示的值(\allowbreak{}None + strings 说明)，不能是 NaN/\allowbreak{}Inf & 67 帧全部通过 pack\_frame；最慢单帧 0.40s；非有限 meta 0 个 \checkmark \\
A & 帧 meta 与内核重算的最大相对差 & 帧的每个读数都由内核函数直接给出，帧层不做第二套计算 & max rel diff =\allowbreak{} 0.000e+00；log10\_lam\_oracle: -1.8；log10\_lam\_lcurve: -3.13333；log10\_lam\_gcv: -2.33333；error\_oracle: 0.00772735 \checkmark \\
A & 同自由度(dof = k\_\allowbreak{}TSVD)下的三方误差与赢家 & 形状不同不等于精度不同：必须在同一条自由度轴上比误差 & matched\_dof =\allowbreak{} 8；tsvd 0.00350，tikhonov 0.00283，landweber 0.09240 \checkmark \\
A & 误差-自由度曲线的形状：TSVD 是离散点、Tikhonov 是连续曲线 & TSVD 的自由度只能取整数 k；Tikhonov 的 lambda 连续 =\allowbreak{}> 自由度连续 & TSVD 曲线 20 点(\allowbreak{}自由度 1..20)；Tikhonov 曲线 61 点(\allowbreak{}自由度 0.01..19.99) \checkmark \\
A & 发散帧与正常帧的掩码零点数 / 因子越界幅度 & omega\_rel > 2 时 |1 - omega sigma\textasciicircum{}2| > 1，滤波因子离开 [0,\allowbreak{}1]，帧里被夹住并标记 & omega\_rel=\allowbreak{}2.5: max|f| =\allowbreak{} 1.918e+05，掩码 0 的个数 1；omega\_rel=\allowbreak{}1.0: 掩码 0 的个数 0 \checkmark \\
A & 发散状态在 strings 里有说明(不是只留一个数) & 占位值三件套：地板/\allowbreak{}夹住的值 + 掩码 + 文案 & omega\_note 提到越界与掩码：True \checkmark \\
A & 两种角度下的零空间维数 / 无噪时的误差地板 & A 是 m x n 时 rank <=\allowbreak{} m；m < n 必有零空间，零空间分量无法从数据恢复 & n\_angles=\allowbreak{}8: null =\allowbreak{} 20，误差 0.1577；n\_angles=\allowbreak{}16: null =\allowbreak{} 0，误差 4.28e-08 \checkmark \\
A & 秩亏与良态两种状态在 strings 里都能读到 & 用户看到的是图，图注必须能解释误差为什么不再下降 & null\_space\_note 提到零空间：True \checkmark \\
A & 帧内解析 / 蒙特卡洛总误差曲线的最大相对差 & 64 个噪声样本的经验均值误差 \textasciitilde{} sigma/\allowbreak{}sqrt(\allowbreak{}64) =\allowbreak{}> 相对差 10\% 量级 & max rel diff =\allowbreak{} 0.099；n\_trials =\allowbreak{} 64 \checkmark \\
A & 大方差端的对数-对数斜率 meta(理论 -4) & lambda >> sigma\_1 时 Var \textasciitilde{} lambda\textasciicircum{}(\allowbreak{}-4) & variance\_slope\_high =\allowbreak{} -3.9993 \checkmark \\
A & 两侧失败方式的读数(lambda 最小 / 最大处的相对误差) & 左端噪声接管、\allowbreak{}右端解被压成 0；两者都应当明显大于最小值 & error\_at\_lam\_min =\allowbreak{} 1.529，error\_at\_lam\_max =\allowbreak{} 1.269，最小误差 0.0199 \checkmark \\
A & 对数轴曲线的最小值(都应当 >= 地板) & 对数轴上 0 画不出来，曲线抬到 1e-17 并说明 & log10\_residual 7.012e-04，error\_curve 7.727e-03，gcv 1.307e-09，curvature 1.000e-17 \checkmark \\
A & Picard 系数落到噪声之下的掩码(0/1 且有 0 也有 1) & 噪声接管之后 |u\_i\textasciicircum{}T b| < ||e||，这一段不能用来判断 Picard 条件 & 掩码和 =\allowbreak{} 15/\allowbreak{}20(\allowbreak{}meta 记 15) \checkmark \\
A & null(无有限值)在 strings 里有专门说明 & meta 里的 null 必须能解释成没测到，而不是 0 & null\_note 存在且含 null：True \checkmark \\
B & lcurve\_\allowbreak{}tradeoff 提供 requiredAux=error & $\le$ 2 & 1 \checkmark \\
B & filter\_\allowbreak{}shapes 提供 requiredAux=error & $\le$ 4 & 1 \checkmark \\
B & image\_\allowbreak{}deconvolution 提供 requiredAux=error & $\le$ 3 & 1 \checkmark \\
B & image\_\allowbreak{}deconvolution 提供 requiredAux=residual & $\le$ 2 & 1 \checkmark \\
B & bias\_\allowbreak{}variance 提供 requiredAux=error & $\le$ 2 & 1 \checkmark \\
B & lcurve\_\allowbreak{}tradeoff 的 theory 条数 & $\le$ 3 & 1 \checkmark \\
B & lcurve\_\allowbreak{}tradeoff 的 pitfalls 长度 & $\le$ 359 & 60 \checkmark \\
B & lcurve\_\allowbreak{}tradeoff 的 pitfalls 含数字 & $\le$ 32 & 1 \checkmark \\
B & lcurve\_\allowbreak{}tradeoff/log10\_\allowbreak{}lam 的 criticalNote 非空 & $\le$ 120 & 1 \checkmark \\
B & filter\_\allowbreak{}shapes 的 theory 条数 & $\le$ 3 & 1 \checkmark \\
B & filter\_\allowbreak{}shapes 的 pitfalls 长度 & $\le$ 307 & 60 \checkmark \\
B & filter\_\allowbreak{}shapes 的 pitfalls 含数字 & $\le$ 10 & 1 \checkmark \\
B & filter\_\allowbreak{}shapes/log10\_\allowbreak{}lam 的 criticalNote 非空 & $\le$ 87 & 1 \checkmark \\
B & filter\_\allowbreak{}shapes/omega\_\allowbreak{}rel 的 criticalNote 非空 & $\le$ 157 & 1 \checkmark \\
B & image\_\allowbreak{}deconvolution 的 theory 条数 & $\le$ 3 & 1 \checkmark \\
B & image\_\allowbreak{}deconvolution 的 pitfalls 长度 & $\le$ 271 & 60 \checkmark \\
B & image\_\allowbreak{}deconvolution 的 pitfalls 含数字 & $\le$ 14 & 1 \checkmark \\
B & image\_\allowbreak{}deconvolution/blur\_\allowbreak{}sigma 的 criticalNote 非空 & $\le$ 95 & 1 \checkmark \\
B & bias\_\allowbreak{}variance 的 theory 条数 & $\le$ 3 & 1 \checkmark \\
B & bias\_\allowbreak{}variance 的 pitfalls 长度 & $\le$ 264 & 60 \checkmark \\
B & bias\_\allowbreak{}variance 的 pitfalls 含数字 & $\le$ 16 & 1 \checkmark \\
B & bias\_\allowbreak{}variance/log10\_\allowbreak{}lam 的 criticalNote 非空 & $\le$ 68 & 1 \checkmark \\
B & lcurve\_\allowbreak{}tradeoff/log10\_\allowbreak{}lam critical >= min & $\le$ 0 & 0 \checkmark \\
B & lcurve\_\allowbreak{}tradeoff/log10\_\allowbreak{}lam critical <= max & $\le$ 0 & 0 \checkmark \\
B & filter\_\allowbreak{}shapes/log10\_\allowbreak{}lam critical >= min & $\le$ 0 & 0 \checkmark \\
B & filter\_\allowbreak{}shapes/log10\_\allowbreak{}lam critical <= max & $\le$ 0 & 0 \checkmark \\
B & filter\_\allowbreak{}shapes/omega\_\allowbreak{}rel critical >= min & $\le$ 0 & 0 \checkmark \\
B & filter\_\allowbreak{}shapes/omega\_\allowbreak{}rel critical <= max & $\le$ 0 & 0 \checkmark \\
B & filter\_\allowbreak{}shapes/omega\_\allowbreak{}rel critical >= min & $\le$ 0 & 0 \checkmark \\
B & filter\_\allowbreak{}shapes/omega\_\allowbreak{}rel critical <= max & $\le$ 0 & 0 \checkmark \\
B & image\_\allowbreak{}deconvolution/blur\_\allowbreak{}sigma critical >= min & $\le$ 0 & 0 \checkmark \\
B & image\_\allowbreak{}deconvolution/blur\_\allowbreak{}sigma critical <= max & $\le$ 0 & 0 \checkmark \\
B & bias\_\allowbreak{}variance/log10\_\allowbreak{}lam critical >= min & $\le$ 0 & 0 \checkmark \\
B & bias\_\allowbreak{}variance/log10\_\allowbreak{}lam critical <= max & $\le$ 0 & 0 \checkmark \\
C & 最贵参数组合下的单帧耗时 & 交互式可视化要求单帧 <=\allowbreak{} 5s(\allowbreak{}预算)，帧层不得有隐藏的 O(\allowbreak{}n\textasciicircum{}3) 重算 & lcurve\_tradeoff 0.01s；filter\_shapes 0.04s；image\_deconvolution 0.50s；bias\_variance 0.31s(\allowbreak{}最慢 0.50s) \checkmark \\
A & 曲率拐点下标(本模块 vs 1.3 lcurve\_\allowbreak{}corner) & 同一离散曲率公式(\allowbreak{}相邻三点中心差分)=\allowbreak{}> 下标必须逐位相同 & noise=\allowbreak{}1e-2: 37 vs 37；noise=\allowbreak{}1e-3: 29 vs 29；noise=\allowbreak{}1e-4: 22 vs 22；noise=\allowbreak{}1e-5: 14 vs 14；noise=\allowbreak{}1e-6: 4 vs 4 \checkmark \\
A & 三个拐点判据的 lambda 极差 / 下标跨度 & 曲率、\allowbreak{}两段直线交点、\allowbreak{}离散二阶差分是三个不同启发式，没有理由给出同一个数 & kappa=\allowbreak{}1e3: lambda 极差 681 倍、\allowbreak{}下标跨度 34(\allowbreak{}曲率 idx 5，两段直线 idx 39，二阶差分 idx 26)；kappa=\allowbreak{}1e4: lambda 极差 126 倍、\allowbreak{}下标跨度 21(\allowbreak{}曲率 idx 20，两段直线 idx 41，二阶差分 idx 32)；kappa=\allowbreak{}1e5: lambda 极差 33 倍、\allowbreak{}下标跨度 13(\allowbreak{}曲率 idx 24，两段直线 idx 31，二阶差分 idx 37)；kappa=\allowbreak{}1e6: lambda 极差 40 倍、\allowbreak{}下标跨度 12(\allowbreak{}曲率 idx 29，两段直线 idx 41，二阶差分 idx 40)；kappa=\allowbreak{}1e7: lambda 极差 63 倍、\allowbreak{}下标跨度 12(\allowbreak{}曲率 idx 33，两段直线 idx 45，二阶差分 idx 42) \checkmark \\
A & 无约束 vs 有约束两段直线拟合：一致组数 / 最大偏离倍数(18 组配置) & 无约束时段越短拟合残差越小 =\allowbreak{}> 切点无规律漂移；约束把切点限制在 [25\%,\allowbreak{} 75\%] 且交点必须落在包围盒内 & 一致 7/\allowbreak{}18 组；其余偏离最大 125.9 倍(\allowbreak{}中位 1.39) \checkmark \\
A & 有约束拐点的 lambda 落在 L 曲线拐弯区间内 & 约束版的交点应当落在曲率极大点附近(\allowbreak{}同一段弯)，偏离量就是这条判据的不确定度 & 有约束 lambda =\allowbreak{} 2.929e-02，曲率判据 lambda =\allowbreak{} 7.356e-04，比值 39.81 \checkmark \\
A & GCV 曲线与 1.3 cross\_\allowbreak{}validation\_\allowbreak{}curve 的最大相对差 & G(\allowbreak{}lambda) =\allowbreak{} ||A x\_lambda - b||\textasciicircum{}2 /\allowbreak{} (\allowbreak{}m - sum f\_i)\textasciicircum{}2，两条实现同一个公式 & max rel diff =\allowbreak{} 3.502e-12；最小值位置 0.00522335 vs 0.00522335 \checkmark \\
A & GCV 不被噪声水平骗：lambda\_\allowbreak{}GCV 随噪声下降而下降 & 噪声越小，能安全保留的模态越多 =\allowbreak{}> 最优 lambda 越小 & lambda\_GCV =\allowbreak{} [0.016 0.005 0.002 0.001 0.   ](\allowbreak{}噪声 1e-2 -> 1e-6) \checkmark \\
A & lambda\_\allowbreak{}GCV 与 GCV 曲线最小点的相对差 & lambda\_GCV =\allowbreak{} argmin G(\allowbreak{}lambda)(\allowbreak{}定义) & |lambda\_GCV - lambda\_argmin| /\allowbreak{} lambda\_GCV =\allowbreak{} 0.000e+00 \checkmark \\
A & 偏差原则处的残差 / 噪声水平 & ||A x\_lambda - b|| =\allowbreak{} tau ||e||(\allowbreak{}tau =\allowbreak{} 1)，交点用对数插值 & lambda\_disc =\allowbreak{} 1.3396e-02，比值 =\allowbreak{} 0.985457 \checkmark \\
A & oracle 位置与误差曲线最小点 & lambda\_oracle =\allowbreak{} argmin ||x\_lambda - x*||(\allowbreak{}定义) & |lambda\_oracle - lambda\_argmin| /\allowbreak{} lambda\_oracle =\allowbreak{} 0.000e+00；最小误差 0.0077 \checkmark \\
A & 四个判据的误差放大倍数(oracle / 偏差原则 / GCV / L 曲线 / 不正则化) & 判据越依赖残差小这种间接信息，越偏向欠正则化的一端 & 1.000(\allowbreak{}oracle)| 1.159(\allowbreak{}disc)| 3.791(\allowbreak{}GCV)| 21.478(\allowbreak{}L)| 12863(\allowbreak{}无) \checkmark \\
A & 网格两端(lambda 最小 / 最大)的相对误差与解范数比 & 左端：无正则化的噪声放大；右端：x\_lambda -> 0 =\allowbreak{}> 误差 -> ||x*|| & 左端 [1.241e-01 1.115e+00 1.040e+01 9.940e+01 9.673e+02]；右端 [0.993 0.993 0.992 0.992 0.992]；||x(\allowbreak{}lambda\_max)||/\allowbreak{}||x(\allowbreak{}lambda\_min)|| =\allowbreak{} [7.50e-03 5.35e-03 8.05e-04 8.78e-05 9.28e-06] \checkmark \\
A & 曲线坐标与逐点独立求解的最大相对差(||x|| 与相对残差) & L 曲线的每个点就是该 lambda 处 Tikhonov 解的范数与残差 & ||x|| 相对差 0.000e+00，相对残差 0.000e+00 \checkmark \\
A & current\_\allowbreak{}solution 与 tikhonov 的一致性(拖动 lambda 的那条路) & 交互拖动 lambda 时前端拿到的是同一个解 & max|dx| =\allowbreak{} 0.000e+00 \checkmark \\
B & 三个拐点下标都在网格内部(kappa=1e3) & $\le$ 0 & 0 \checkmark \\
B & 三个拐点下标都在网格内部(kappa=1e5) & $\le$ 0 & 0 \checkmark \\
B & 三个拐点下标都在网格内部(kappa=1e7) & $\le$ 0 & 0 \checkmark \\
B & 有效判据读数全有限 & 0 个非有限值 & 0 \checkmark \\
B & 无噪时偏差原则没有交点 -> NaN(不是 0) & 1 & 1 \checkmark \\
B & ||x\_\allowbreak{}lambda|| 随 lambda 不增 & nonincreasing & 违反 0 处 \checkmark \\
B & 相对残差随 lambda 不减 & nondecreasing & 违反 0 处 \checkmark \\
B & 离散曲率全有限 & 0 个非有限值 & 0 \checkmark \\
B & 网格左端 = sigma\_\allowbreak{}n/10 & $1.000\times 10^{-7}$ & $1.000\times 10^{-7}$ \checkmark \\
B & 网格右端 = 10 sigma\_\allowbreak{}1 & 10 & 10 \checkmark \\
B & 网格单调递增 & increasing & 违反 0 处 \checkmark \\
C & L 曲线代价降到 10 倍以下的噪声水平(二分) & 最优 lambda 随噪声下降，而 L 曲线拐点滞后 =\allowbreak{}> 噪声越小拐点越差 & 阈值 =\allowbreak{} 4.1727e-03(\allowbreak{}对数 2.38 个数量级)；噪声 1e-2..1e-6 时的代价 [  6.5  21.5  30.  199.1 211.6] \checkmark \\
C & lambda\_\allowbreak{}GCV 对噪声水平的幂律指数 & 几何谱 + GCV：lambda\_GCV \textasciitilde{} noise\textasciicircum{}(\allowbreak{}1/\allowbreak{}2)(\allowbreak{}实测指数略低于 1/\allowbreak{}2) & p =\allowbreak{} 0.4667(\allowbreak{}R\textasciicircum{}2 =\allowbreak{} 0.997557)；lambda\_GCV =\allowbreak{} [0.02 0.   0.   0.   0.  ] \checkmark \\
C & lambda\_\allowbreak{}oracle 对噪声水平的幂律指数(对照) & oracle 的 lambda 也随噪声下降；GCV 的指数与它的差就是判据的滞后量 & p =\allowbreak{} 0.4800(\allowbreak{}GCV 比它小 0.013) \checkmark \\
C & 网格左端误差对 kappa 的幂律指数 & 不做正则化时误差 \textasciitilde{} ||e||/\allowbreak{}(\allowbreak{}sigma\_n ||x*||) \textasciitilde{} kappa & p =\allowbreak{} 0.9733(\allowbreak{}R\textasciicircum{}2 =\allowbreak{} 0.999950)；误差 [1.241e-01 1.115e+00 1.040e+01 9.940e+01 9.673e+02] \checkmark \\
A & 滤波因子与闭式的最大偏差 / 半功率点 & f\_i =\allowbreak{} s\_i\textasciicircum{}2/\allowbreak{}(\allowbreak{}s\_i\textasciicircum{}2+lambda\textasciicircum{}2)；s\_i =\allowbreak{} lambda 处 f =\allowbreak{} 1/\allowbreak{}2 & max|df| =\allowbreak{} 0.000e+00；f(\allowbreak{}s=\allowbreak{}lambda) =\allowbreak{} 0.500000000000 \checkmark \\
A & lambda=0 与最小二乘解的一致性 / lambda 极大时的压缩比 & lambda=\allowbreak{}0 时 Tikhonov =\allowbreak{} 最小二乘；lambda >> sigma\_1 时 ||x|| \textasciitilde{} ||x\_ls|| sigma\_1\textasciicircum{}2/\allowbreak{}lambda\textasciicircum{}2 & rel diff =\allowbreak{} 2.819e-11；||x\_huge||/\allowbreak{}||x\_ls|| =\allowbreak{} 8.796e-15 \checkmark \\
A & 谱公式与 LU 正规方程解的最大相对差 & <=\allowbreak{} 1e-10(\allowbreak{}lambda 不太小时两条路线同一个解；kappa=\allowbreak{}1e6 下实测 2e-12) & lambda=\allowbreak{}0.01: 1.97e-12；lambda=\allowbreak{}0.1: 1.91e-14；lambda=\allowbreak{}1: 7.36e-16；lambda=\allowbreak{}0.01: 1.97e-12 \checkmark \\
A & 正规方程路线在 lambda -> 0 时的退化(kappa 平方的代价) & 正规方程的条件数是 kappa(\allowbreak{}A)\textasciicircum{}2；lambda 小到正则项压不住时 LU 解开始丢位 & lambda=\allowbreak{}1e-2: 1.97e-12；lambda=\allowbreak{}1e-6: 3.19e-06(\allowbreak{}劣化 1614925 倍；kappa(\allowbreak{}A)\textasciicircum{}2 =\allowbreak{} 1.000e+12) \checkmark \\
A & TSVD 与 1.3 truncated\_\allowbreak{}solution 的最大逐点差 / 误差曲线 & k =\allowbreak{} n 时退化为最小二乘解；误差在 k* 处最小、\allowbreak{}之后被噪声抬高 & max|dx| =\allowbreak{} 5.551e-17；k* =\allowbreak{} 7(\allowbreak{}err 0.0163)，k =\allowbreak{} n 时 err =\allowbreak{} 100.1523(\allowbreak{}放大 6141.3 倍) \checkmark \\
A & 谱形式与矩阵迭代的最大相对差 / 发散线另一侧的表现 & 两者同解(\allowbreak{}<=\allowbreak{} 1e-9)；omega >=\allowbreak{} 2/\allowbreak{}sigma\_1\textasciicircum{}2 时迭代发散、\allowbreak{}滤波因子离开 [0,\allowbreak{}1] & omega\_rel=\allowbreak{}1,\allowbreak{}k=\allowbreak{}30: 6.63e-14(\allowbreak{}stop=\allowbreak{}n\_iter)；omega\_rel=\allowbreak{}0.5,\allowbreak{}k=\allowbreak{}60: 2.61e-13(\allowbreak{}stop=\allowbreak{}n\_iter)；omega\_rel=\allowbreak{}1.5,\allowbreak{}k=\allowbreak{}12: 3.18e-14(\allowbreak{}stop=\allowbreak{}n\_iter)；omega\_rel=\allowbreak{}2.2: stop=\allowbreak{}diverged，max f =\allowbreak{} 6.859e+15 \checkmark \\
A & 两点 TV 闭式解的最大逐点偏差 & x =\allowbreak{} (\allowbreak{}m - soft(\allowbreak{}d,\allowbreak{} lambda),\allowbreak{} m + soft(\allowbreak{}d,\allowbreak{} lambda))，m =\allowbreak{} mean(\allowbreak{}b)，d =\allowbreak{} (\allowbreak{}b\_2-b\_1)/\allowbreak{}2 & lambda=\allowbreak{}0: max|dx|=\allowbreak{}0.00e+00；lambda=\allowbreak{}0.05: max|dx|=\allowbreak{}1.11e-16；lambda=\allowbreak{}0.3: max|dx|=\allowbreak{}1.11e-16；lambda=\allowbreak{}0.75: max|dx|=\allowbreak{}5.55e-17；lambda=\allowbreak{}2: max|dx|=\allowbreak{}1.11e-16 \checkmark \\
A & lambda >= |d| 时 TV 解是否退化成常数 & TV 罚项在差值上的次梯度包含 0 时，两点被合并(\allowbreak{}台阶先验) & lambda=\allowbreak{}1.0: x =\allowbreak{} (\allowbreak{}0.250000000,\allowbreak{} 0.250000000)，差 =\allowbreak{} 1.11e-16 \checkmark \\
A & lambda=0 时 TV 探针是否退化为恒等(正对照，kappa=1e2) & lambda=\allowbreak{}0 时 TV 项不起作用，x =\allowbreak{} A\textasciicircum{}+ b，探针读数应为 1 & lambda=\allowbreak{}0 的探针读数最小偏差 =\allowbreak{} 5.45e-14(\allowbreak{}2000 步) \checkmark \\
A & TV 探针的依赖性 / 非对角泄漏 & TV 非线性且非对角：换探针读数会变，泄漏不为 0(\allowbreak{}它不是滤波因子) & max|f\_mixed - f\_single| =\allowbreak{} 0.959；max leakage =\allowbreak{} 0.144 \checkmark \\
A & 蒙特卡洛总误差落在解析值的 4 倍标准误内 & 64 个样本的经验均值误差 \textasciitilde{} sigma/\allowbreak{}sqrt(\allowbreak{}64)；逐点比较要用标准误刻度 & 越界点 0/\allowbreak{}41；最大相对差 mse 0.099，variance 0.118 \checkmark \\
A & 偏差\textasciicircum{}2 + 方差 = 总误差(恒等式) & 两项按定义相加就是 E||x\_lambda - x*||\textasciicircum{}2 & max|bias2+variance-mse| =\allowbreak{} 0.000e+00 \checkmark \\
A & lambda 最小时实测偏差\textasciicircum{}2 与解析偏差\textasciicircum{}2 的绝对差(以 mse\_\allowbreak{}min 为单位) & 经验均值 E\_hat[x] 的抽样误差是 sigma/\allowbreak{}sqrt(\allowbreak{}N)，它进入偏差\textasciicircum{}2 后不可忽略 & bias2\_analytic =\allowbreak{} 6.407e-09，bias2\_mc =\allowbreak{} 4.354e-02，绝对差 =\allowbreak{} 4.354e-02 =\allowbreak{} 106.5 x mse\_min \checkmark \\
A & 噪声对齐最小奇异方向时的相对误差与解析值 & ||A\textasciicircum{}+ e|| =\allowbreak{} ||e||/\allowbreak{}sigma\_n(\allowbreak{}e 沿 u\_n)，相对误差 =\allowbreak{} ||e||/\allowbreak{}(\allowbreak{}sigma\_n ||x*||) & 实测 900.363172 /\allowbreak{} 解析 900.363172(\allowbreak{}相对差 7.92e-12) \checkmark \\
A & kappa=1e6、噪声 1e-3 时的放大倍数(最坏 / 均方根) & 最坏 ||e||sqrt(\allowbreak{}sum 1/\allowbreak{}s\_i\textasciicircum{}2)/\allowbreak{}||x*||，均方根除以 sqrt(\allowbreak{}m) & 最坏 1028.4 倍，均方根 181.8 倍(\allowbreak{}两者相差 5.7 倍) \checkmark \\
A & 无噪 Picard 斜率 / 比值斜率 / 相关系数 & alpha\_i =\allowbreak{} sigma\_i\textasciicircum{}decay(\allowbreak{}decay=\allowbreak{}2)=\allowbreak{}> 斜率 =\allowbreak{} 2，比值斜率 =\allowbreak{} 1，log-log 完全共线 & slope =\allowbreak{} 2.000000，ratio slope =\allowbreak{} 1.000000，corr =\allowbreak{} 1.000000 \checkmark \\
A & 有噪时整段斜率 vs 信号段斜率 / 被噪声接管的模态数 & 噪声接管后 |u\_i\textasciicircum{}T b| 不再随 sigma\_i 下降，整段斜率必然变平；信号段仍满足理论 & 整段 0.5641，信号段 1.9972(\allowbreak{}5 个模态)，低于噪声的模态 15(\allowbreak{}按 sigma\_res 计 10) \checkmark \\
A & Picard 不成立时：oracle 位置与正则化的收益 & |u\_i\textasciicircum{}T b|/\allowbreak{}sigma\_i 随 i 增长时，最小二乘解在噪声下必然崩，正则化只能截掉一部分 & oracle lambda =\allowbreak{} 1.000e-07(\allowbreak{}网格左端 1.000e-07)，误差 0.116，正则化收益 =\allowbreak{} 1.000 倍 \checkmark \\
A & Picard 成立时的对照：最优 lambda 处误差与不正则化的比 & 同一算子、\allowbreak{}同一噪声，Picard 成立时正则化把误差从 1e4 量级压回 1e-2 量级 & 不正则化误差 99.399，oracle 处 0.0077，比值 12863 倍 \checkmark \\
A & 二维 Tikhonov 解的正规方程残差(相对) & (\allowbreak{}B\textasciicircum{}T B) X (\allowbreak{}B\textasciicircum{}T B) + lambda\textasciicircum{}2 X =\allowbreak{} B\textasciicircum{}T Y B 是二维 Tikhonov 的定义方程 & max|残差|/\allowbreak{}max|B\textasciicircum{}T Y B| =\allowbreak{} 5.130e-15 \checkmark \\
A & 二维条件数 = 一维条件数的平方 & 可分离算子的奇异值是 s\_i s\_j，因此 kappa\_2D =\allowbreak{} kappa\_1D\textasciicircum{}2 & kappa\_1D =\allowbreak{} 5911 -> kappa\_2D =\allowbreak{} 3.494e+07(\allowbreak{}s\_1 s\_1 =\allowbreak{} 1，s\_n s\_n =\allowbreak{} 2.862e-08) \checkmark \\
A & 停滞 SVD 与体检后 SVD 的相对误差(对 LAPACK) & 单边 Jacobi 在重奇异值矩阵上可能停滞；factorize 用 converged + 重建 + 正交性三道体检回退 & jacobi: converged=\allowbreak{}False，相对误差 2221.2 倍；factorize -> kogbetliantz(\allowbreak{}fallback=\allowbreak{}1)，相对误差 1.88e-13 \checkmark \\
A & CT 投影矩阵的行和 / 伴随恒等式 / 两种角度下的秩 & 射线驱动的线积分归一化；A 的秩由角度数决定，角度不足时零空间非平凡 & 每行行和与该角度弦长的最大偏差 5.329e-15(\allowbreak{}8 个角度 x 10 条射线)；|<Ax,\allowbreak{}y>-<x,\allowbreak{}A\textasciicircum{}T y>| =\allowbreak{} 5.773e-15；8 角度秩 80/\allowbreak{}100，14 角度秩 100/\allowbreak{}100，kappa 从 1.97e+05 降到 1.78e+01 \checkmark \\
B & ||x\_\allowbreak{}lambda|| 随 lambda 不增 & nonincreasing & 违反 0 处 \checkmark \\
B & 相对残差随 lambda 不减 & nondecreasing & 违反 0 处 \checkmark \\
B & 有效自由度随 lambda 不增 & nonincreasing & 违反 0 处 \checkmark \\
B & TSVD 残差随 k 不增 & nonincreasing & 违反 0 处 \checkmark \\
B & TSVD 滤波因子只有 0/1 & 0 & 0 \checkmark \\
B & Landweber 上界(omega\_\allowbreak{}rel=0.1) & $\le$ 1 & 1 \checkmark \\
B & Landweber 下界(omega\_\allowbreak{}rel=0.1) & $\le$ 0 & 0 \checkmark \\
B & Landweber 上界(omega\_\allowbreak{}rel=0.5) & $\le$ 1 & 1 \checkmark \\
B & Landweber 下界(omega\_\allowbreak{}rel=0.5) & $\le$ 0 & 0 \checkmark \\
B & Landweber 上界(omega\_\allowbreak{}rel=1.0) & $\le$ 1 & 1 \checkmark \\
B & Landweber 下界(omega\_\allowbreak{}rel=1.0) & $\le$ 0 & 0 \checkmark \\
B & Landweber 上界(omega\_\allowbreak{}rel=1.9) & $\le$ 1 & 1 \checkmark \\
B & Landweber 下界(omega\_\allowbreak{}rel=1.9) & $\le$ 0 & 0 \checkmark \\
B & 自由度上界 & $\le$ 20 & 10 \checkmark \\
B & 最优步之前的误差不增 & nonincreasing & 违反 0 处 \checkmark \\
B & 最优步处误差小于末步(半收敛) & $\le$ 0.035007 & 0.035007 \checkmark \\
B & 误差历史全有限 & 0 个非有限值 & 0 \checkmark \\
B & TV 目标函数不超过候选点 & $\le$ 0.0615792 & 0.0337634 \checkmark \\
B & 幻影只有四档灰度(含外围 1.0) & 0 & 0 \checkmark \\
B & 模糊矩阵行和为 1 & 0 & 0 \checkmark \\
B & 模糊矩阵对称 & 0 & 0 \checkmark \\
C & 矩阵迭代实测的发散阈值(400 步预算) & 2 & 2.07152 \checkmark \\
C & 谱判据给出的发散阈值(= 2/sigma\_\allowbreak{}1\textasciicircum{}2，精确) & 2 & 2 \checkmark \\
C & 方差尾部的幂律指数(lambda 从 1e5 降到 10，25 个点) & lambda >> sigma\_1 时 f\_i \textasciitilde{} s\_i\textasciicircum{}2/\allowbreak{}lambda\textasciicircum{}2 =\allowbreak{}> Var \textasciitilde{} lambda\textasciicircum{}(\allowbreak{}-4)，拟合区间 lambda in [10 sigma\_1,\allowbreak{} 1e5] & p =\allowbreak{} -3.9994(\allowbreak{}窗口 1e+01..1e+05，R\textasciicircum{}2 =\allowbreak{} 1.000000) \checkmark \\
C & 最坏放大倍数对 kappa 的幂律指数 & 最坏放大 \textasciitilde{} ||e|| sqrt(\allowbreak{}sum 1/\allowbreak{}s\_i\textasciicircum{}2)/\allowbreak{}||x*||；几何谱下 sum 1/\allowbreak{}s\_i\textasciicircum{}2 由 1/\allowbreak{}sigma\_n\textasciicircum{}2 主导 =\allowbreak{}> 指数 1 & p =\allowbreak{} 0.9780(\allowbreak{}R\textasciicircum{}2 =\allowbreak{} 0.999942)；放大倍数 [1.3e-01 1.1e+00 1.1e+01 1.0e+02 1.0e+03] \checkmark \\
C & 放大倍数与主项 1/sigma\_\allowbreak{}n 的比(多留的模态贡献多少) & 最坏放大 =\allowbreak{} ||e||/\allowbreak{}sigma\_n * sqrt(\allowbreak{}1 + sigma\_n\textasciicircum{}2 sum\_\{i<n\} 1/\allowbreak{}s\_i\textasciicircum{}2)，第二项随 kappa 递减 & 比值 [1.6134 1.3912 1.2693 1.1932 1.1423](\allowbreak{}kappa=\allowbreak{}1e2 高出 61\%，kappa=\allowbreak{}1e6 只高出 14\%) \checkmark \\
\end{longtable}
